\chapter{Asking Millionaires Who They're Voting For President}

This chapter follows a School of Hard Knocks interview, curated by LazyingArt LLC, in which the host asks two questions at once. One question is practical: how was wealth actually built, lost, protected, or rebuilt? The other is political: how do wealthy operators say they will vote in the 2024 presidential election? The lecture begins by spending outcomes before explanations, then resets in Austin, Texas and moves interview by interview through negotiation, marriage and capital preservation, prediction markets, bankruptcy, trading, restaurant risk, and the distinction between saving and investing. We will keep close to that order, because the force of the lecture lies in the way each new anecdote sharpens the previous one.

\section{Cold Open, Austin, and the Two Questions}

The opening is a splice of consequences before causes. We first hear the end-state fragments: a lawyer who used to be a millionaire, rough yearly figures in the millions, and the election question itself. Only then does the host formalize the experiment. We are in Austin, Texas, and the city is introduced as a dense local field of wealth:
\begin{equation}
N_{\text{millionaires}} \approx 40{,}000,\qquad
N_{\text{billionaires}} > 10.
\end{equation}
That count is not yet an analysis; it is a stage-setting device. It tells us why the host expects this street format to produce not merely opinions, but operating doctrines. The chapter therefore proceeds under a double discipline:
\begin{itemize}
\item for each speaker, identify the wealth mechanism or wealth failure that is being described;
\item then note, without overpromoting it to a theorem, the political preference that the host extracts.
\end{itemize}

This order matters. The mathematically serious spine of the lecture is not electoral ideology. It is leverage, realization, timing, risk, and the survival of skill after capital has disappeared.

\section{Buy Dirt, Fail First, and Negotiate Without Desperation}

The first full interview is with a New York builder-developer. The host begins with broke-versus-wealthy, then line of business, and only then lets the doctrine emerge. The speaker's first compressed principle is ancestral and concrete: buy dirt. He works both commercial and residential real estate, points not to a single trophy transaction but to streets named after his children, and shifts quickly from pride to method.

\paragraph{Anecdote.}
The speaker says the best thing he did was not merely a large deal, but a deal that left a permanent mark: streets named after family. The anecdote prevents the section from becoming a sterile treatise on capital allocation. In this lecture, wealth is repeatedly measured not only by gross figures, but by what survives.

\paragraph{Claim.}
Asked for the secret to creating wealth in real estate, the answer is failure. One must fail, then get hungry. The claim is not romanticized. It is tied to a generation difference: earlier real-estate conditions made certain acquisitions easier than they are now.

\paragraph{Mechanism.}
The older mechanism is described in the loose but recognizable phrase ``buy stuff subject to,'' with the further claim that it could gain value even before closing. We should keep that cautiously schematic:
\begin{equation}
\text{buy subject to} \;\Rightarrow\; \text{asset appreciates before closing}.
\end{equation}
The lecture does not supply a legal derivation of ``subject to.'' What matters is the speaker's contrast between an earlier market with embedded lift and a present market that feels riskier.

The section then narrows to negotiation. The host asks for advice at the table with a ``heavy hitter,'' and the speaker answers with a doctrine of anti-desperation. Do not impulse the deal. Let the other side feel that you do not really want it. Do not be desperate. The host later compresses that into the reusable rule:
\begin{equation}
\text{desperation} \Rightarrow \text{loss of leverage},
\qquad
\text{willingness to walk away} \Rightarrow \text{bargaining power}.
\end{equation}

The speaker also gives a numerical example in the language of bank financing:
\begin{equation}
r_{\text{ask}}=6\%,\qquad r_{\text{counter}}=4\%.
\end{equation}
His point is not that every negotiation yields precisely this cut. The point is that the first quote is not sacred, and that a credible exit can reopen a deal that seemed settled.

\subsection{Question \& Answer}

Why does willingness to walk away create leverage?

Because the walk-away option changes the direction of pressure. The speaker's arithmetic is simple but revealing. If the first ask is \(6\%\) and the counter is \(4\%\), then the negotiated movement is
\begin{align}
\Delta r &= r_{\text{ask}}-r_{\text{counter}} = 2\%,\\
\frac{\Delta r}{r_{\text{ask}}} &= \frac{2}{6} \approx 0.33.
\end{align}
So the anecdotal reduction is roughly one third of the original ask. We should not universalize that number. But the structure is useful: the person who must close today cannot force the rate to move, while the person who can leave may turn the other side's fear of losing the deal into a concession. In that limited and commercial sense, the host's recap is mathematically apt: the most desperate person has the smallest feasible set of outcomes.

The interview ends, as the lecture keeps requiring, with a political answer. The speaker says he will vote for Trump, not because the lecture has proven a political theorem, but because the business doctrine has been extracted and the format demands the second question.

\section{Marriage, Paper Wealth, and the Arithmetic of Not Selling}

The next interview changes register. The host begins with the millionaire question, and the speaker answers by denying current millionaire status while immediately invoking loss through divorce. From there the lecture broadens from negotiation at the table to the economics of household structure, prenuptial agreements, and the difference between visible wealth and realized wealth.

\paragraph{Anecdote.}
The speaker is a lawyer. He says several million dollars were lost in divorce, that his former spouse retained the house, and that he later chose separation rather than another divorce because he did not want retirement assets divided away. The host then reframes the whole exchange in a memorable sentence: marriage is a business.

\paragraph{Claim.}
What matters for the notes is not the speaker's universalization about divorce law, but the capital-preservation claim beneath it: if one partner is the main earner and wealth is being built, household structure is economically consequential. In the speaker's own language, a prenuptial agreement belongs to wealth-building.

\paragraph{Mechanism.}
The lecture sharpens when the host asks for the most money made in a single year. The speaker explicitly corrects the category: it was paper money, not realized money. That distinction is the real mathematical payload of the segment.

A convenient formulation is:
\begin{equation}
V_{\text{paper}} \neq V_{\text{realized}}
\qquad
\text{when sale constraints bind.}
\end{equation}
The transcript then supplies the numbers:
\begin{align}
V_{\text{peak}} &\approx \$17.5\,\text{million},&
\Delta t_{\text{lockup}} &= 6\ \text{months},\\
V_{\text{post-crash}} &\approx \$750{,}000,&
V_{\text{sale}} &\approx \$350{,}000.
\end{align}
In compressed narrative form, the episode is
\begin{equation}
\$17.5\text{M} \longrightarrow \$750{,}000 \longrightarrow \$350{,}000.
\end{equation}
The intervening mechanism is precise enough for notes: stock options in dot-com-era IPO companies, followed by an insider lockup, followed by a crash during the lockup, followed by sale only after severe deterioration.

\subsection{Question \& Answer}

When is wealth only paper wealth?

When valuation outruns realizability. In this segment, the shares are marked at a large number, but the owner cannot liquidate on his own timetable. The lockup turns market price into a constrained observation rather than an available cash position.

A short worked calculation makes the point vivid. From the transcript values,
\begin{align}
\frac{V_{\text{post-crash}}}{V_{\text{peak}}} &\approx \frac{0.75}{17.5} \approx 0.043,\\
\frac{V_{\text{sale}}}{V_{\text{peak}}} &\approx \frac{0.35}{17.5} \approx 0.020.
\end{align}
So the realized exit is about \(2\%\) of the paper peak. Equivalently,
\begin{equation}
1-\frac{V_{\text{sale}}}{V_{\text{peak}}} \approx 0.98.
\end{equation}
About \(98\%\) of the peak paper value disappears before realization. This is why the speaker's line, ``sell too soon,'' matters. It is not a theorem about every asset. It is an attack on the fantasy that quoted wealth is already secured wealth.

The host then asks the election question. The speaker refuses to discuss politics in public but gives a directional preference against Kamala Harris. That answer is kept as a speaker-attributed conclusion, not as the analytic center of the section.

\section{Election Markets as a Probability Interruption}

At this point the lecture detours into a sponsor-style block. This interruption should not be mistaken for the settled doctrine of the chapter, but it does something mathematically useful: it translates political speech into priced probabilities. The host says that the general election is trading near even, and the interviewee adds that only a narrow margin is at issue:
\begin{equation}
p_{\text{general}} \approx 50\text{--}50,\qquad
\delta_{\text{race}} \approx 1\%\text{--}2\%.
\end{equation}
The block then adds two more numerical claims:
\begin{equation}
V_{\text{platform}} \approx \$1\,\text{billion},
\qquad
D=\$100 \Longrightarrow C=\$20.
\end{equation}
These are host-side platform claims, not guest doctrine. Still, they matter because the lecture briefly redefines the election as a probability market rather than a polling conversation.

The visible evidence for this shift is the dashboard of card-based markets.

\begin{figure}[t]
\centering
\includegraphics[width=0.95\textwidth]{lecture_61_figure_02.png}
\caption{Prediction-market dashboard with visible state-level presidential odds. The subtitle mentions swing states such as Georgia or Arizona, but the readable cards shown here are California, Iowa, New Mexico, Washington, D.C., and Indiana.}
\end{figure}

The screenshot is better read as a market layout than as a map. The natural notation is
\begin{equation}
p_{s,c}=\Pr(\text{candidate }c\text{ wins state }s).
\end{equation}
Using only the clearly legible cards, we can reconstruct the visible state odds as
\begin{align}
p_{\text{CA,Harris}} &\approx 0.98, & p_{\text{CA,Trump}} &\approx 0.02,\\
p_{\text{IA,Trump}} &\approx 0.95, & p_{\text{IA,Harris}} &\approx 0.05,\\
p_{\text{NM,Harris}} &\approx 0.89, & p_{\text{NM,Trump}} &\approx 0.11,\\
p_{\text{DC,Harris}} &\approx 0.98, & p_{\text{DC,Trump}} &\approx 0.02,\\
p_{\text{IN,Trump}} &\approx 0.97, & p_{\text{IN,Harris}} &\approx 0.03.
\end{align}
A side market on the same dashboard also shows that the platform is not purely political:
\begin{equation}
\Pr\!\bigl(\text{RT score of ``Venom: The Last Dance''}>30\bigr)\approx 0.91,
\qquad
\Pr\!\bigl(\text{RT score}>45\bigr)\approx 0.14.
\end{equation}

The important point is structural. The host has spent the earlier interviews on business doctrine and personal loss; here he temporarily inserts a new language in which uncertainty is priced and distributed across many parallel cards. The chapter should preserve that interruption, because it is the lecture's one clean moment of overt probabilistic notation.

\section{Bankruptcy, Sweat Equity, and the Founder Who Must Be Everyone}

The lecture then returns from priced beliefs to lived business history. A new speaker says he became a millionaire by \(27\), filed bankruptcy by \(31\), and took about ten years to rebuild. The movement is exact enough to warrant a short timeline:
\begin{align}
a_{\text{startup}} &= 24,\qquad
a_{\text{millionaire}} = 27,\qquad
a_{\text{bankruptcy}} = 31,\\
R_{\text{excavation}} &\approx \$8\,\text{million/year at } a\approx 28\text{--}29,\qquad
T_{\text{rebuild}} \approx 10\ \text{years}.
\end{align}

\paragraph{Anecdote.}
The business was excavation. The speaker liked being outside, liked building something from raw ground, and says he tried other trades before discovering the one that fit. The host's phrase ``riches are in the niches'' is not the speaker's theorem, but the interview makes room for it.

\paragraph{Claim.}
The central claim is that nobody is coming to save the founder. There is no owner above the owner, no vice-president above the vice-president. The business must be made to happen by the same person who is trying to learn how it works.

\paragraph{Mechanism.}
The startup mechanism is described as sweat equity. Cash capital is not denied, but it is not the center of the story. The speaker worked equipment by day and went to school by night to learn project management, estimating, and contract negotiation. A cautious reconstruction is
\begin{equation}
K_0^{\text{effective}} \approx K_0^{\text{labor}} + K_0^{\text{skills}},
\end{equation}
with the understanding that this is a note-writer's compression, not a line spoken verbatim.

The organizational compression is equally important:
\begin{equation}
\text{founder}
=
\text{owner}
+
\text{VP}
+
\text{operator}
+
\text{estimator}
+
\text{negotiator}.
\end{equation}
Again, this is not a formal organization chart. It is the lecture's way of saying that early entrepreneurial competence is multiplex.

\subsection{Question \& Answer}

What survives bankruptcy when the money does not?

The lecture's answer is: skills, habits, and the capacity to reassemble a business. The host's recap after this segment is unusually sharp. You can take away the businesses, the money, and the accumulated position, but not necessarily the skill set or the mindset that built them.

We can write the point schematically as
\begin{equation}
\text{capital stock collapses} \;\centernot\Rightarrow\; \text{operating capacity collapses}.
\end{equation}
In the language of the interview, the proof is biographical rather than formal: first wealth, then collapse, then a ten-year rebuild. The lecture is not promising inevitable recovery. It is distinguishing between balance-sheet loss and the destruction of the person who learned to build.

The political answer here is emphatic and again speaker-attributed: the entrepreneur says he is voting for Trump and expects a landslide. The host then recaps the surprising scorecard for Austin before moving on.

\section{Restaurants, Trading, and the Save-versus-Invest Split}

The final major interview is with a retired finance professional who also owned restaurants. This section has a quieter tone, but it gives the lecture its broadest money-language.

\paragraph{Anecdote.}
The speaker says he owned three restaurant locations and had been a bond trader. The host asks about restaurant margins and receives a hard binary description: either enormous success or failure. The transcript gives us the scale markers
\begin{equation}
N_{\text{locations}} = 3,\qquad
R_{\max} \approx \$50\,\text{million}.
\end{equation}
The restaurant risk can be compressed as
\begin{equation}
\text{restaurant outcome} \in \{\text{hugely successful},\ \text{out of business}\}.
\end{equation}

\paragraph{Claim.}
The speaker says banks do not like lending to restaurants and describes the business as privately financed. He then pivots, as a former trader, to a general warning: anyone who says he never lost money is lying. Loss belongs to the education.

\paragraph{Mechanism.}
The business lesson is not that loss is pleasant; it is that downside exposure forces realism. Restaurant margins do it through operational fragility. Trading does it through mark-to-market pain. In both settings, one learns by mispricing risk and paying for that error.

\subsection{Question \& Answer}

Why do restaurants and traders both force respect for loss?

Because both place a decision-maker inside a tight feedback loop. In the restaurant case, costs, margins, and financing pressure make bad assumptions visible quickly. In trading, price movement does the same thing. That is why the speaker treats loss not as an exception to business knowledge, but as one of its sources.

The section then reaches its most portable aphorism. Asked what separates the middle class from wealthy people, the speaker says:
\begin{equation}
\text{middle class} \to \text{saves},
\qquad
\text{upper class} \to \text{invests}.
\end{equation}
A minimal arithmetic rendering, useful for notes but not spoken verbatim, is
\begin{align}
W_{t+1}^{\text{save}} &= W_t + s_t,\\
W_{t+1}^{\text{invest}} &= W_t + s_t + r_t W_t.
\end{align}
The first rule preserves and accumulates principal through savings flow alone. The second adds exposure to return on existing wealth. This is precisely the distinction the lecture has been assembling all along. Real estate leveraged timing and negotiation. Dot-com stock options exposed the difference between quoted value and captured value. Prediction markets priced uncertainty. The excavation story treated skill as recoverable capital. The final speaker now compresses all that into one class distinction.

He refuses a specific investment sector, answering ``whatever,'' and then gives a Peter Lynch-style observational rule: invest by walking around. If the products in ordinary use around you are obvious, the investment universe need not be mysterious. The lecture ends this final substantive segment with a softened political answer: split household, party over individual, Republican affiliation, but uncertainty about the election outcome itself.

\section{Summary}

The chapter opens with outcomes before causes and then proceeds through four distinct money problems. The real-estate developer gives a first theorem of leverage: never negotiate from desperation, and never treat the opening rate as final. The lawyer converts wealth talk into a harder arithmetic: paper wealth under lockup can collapse before it is realized. The sponsor interruption briefly recasts politics as a field of priced probabilities and visible state-level market cards. The excavation entrepreneur insists that bankruptcy can erase capital without erasing the operating person. The final finance-and-restaurant speaker finishes the line of thought by making loss educational and by drawing the sharp distinction between saving and investing.

That is the real through-line of the lecture. Wealth here is not a static possession. It is a sequence of negotiations, constraints, exposures, losses, rebuilds, and update rules. The votes are part of the format. The business logic is the deeper subject.